\subsection{Zásadní rozdíly v přístupu}
\label{subsec:static-vs-llm-differences}

Klíčové rozdíly mezi statickými analyzátory a LLM lze shrnout do čtyř kategorii: determinismus výstupu, hloubka formální analýzy, práce s kontextem a srozumitelnost výstupů.

\subsubsection*{Determinismus vs pravděpodobnost.}
\label{subsubsec:determinism-vs-probability}

Statický analyzátor je deterministický systém, pro identický vstup vždy produkuje identický výstup.
To umožňuje přesnou reprodukovatelnost výsledků a snadné testování samotného nástroje.

LLM je pravděpodobnostní systém, kde výstup závisí na parametrech inference (zejména teplotě) a není plně reprodukovatelný.
Při nízké teplotě se modely výstupu přibližují, ale absolutní reprodukovatelnost není garantována.
Pro systém poskytující automatizovanou zpětnou vazbu studentům to znamená, že výsledky analýzy pro dvě identická odevzdání se mohou v detailech lišit.

\subsubsection*{Formální analýza vs heuristická interpretace.}
\label{subsubsec:formal-vs-heuristic}

Statický analyzátor provádí formálně definovanou analýzu nad přesnou strukturální reprezentací programu.
Může dokazovat (v rámci svého modelu) přítomnost nebo absenci určitého vzoru.

LLM provádí heuristickou interpretaci, jeho závěry jsou odhady odvozené ze statistických vzorů, nikoli z formálního průchodu strukturou programu.
Z toho vyplývá, že LLM může chybět detekce vzoru, který statický analyzátor najde spolehlivě, a naopak tak LLM může upozornit na problém, který statická analýza vůbec nedokáže formalizovat.

\subsubsection*{Kontext zadání a záměr.}
\label{subsubsec:context}

Toto je oblast, kde LLM přináší základně odlišnou hodnotu.
Statický analyzátor nemá přístup k žádnému kontextu vně zdrojového kódu.
Neví, jakou úlohu program řeší, jaké jsou požadavky na výkonnost nebo jaká jsou omezení implementace.
LLM, kterému je v promptu poskytnuto zadání úlohy, dokáže posoudit soulad implementace se zadáním, efektivitu zvoleného algoritmu, vhodnost použitých datových struktur i správné ošetření hraničních vstupů.
Tato schopnost je pro hodnocení studentských prací velmi cenná, jde přesně o tu vrstvu, která vyžaduje lidský úsudek a nejvíce zatěžuje čas vyučujícího.

\subsubsection*{Forma výstupů a srozumitelnost.}
\label{subsubsec:output-format}

Statické analyzátory typicky produkují stručná technická hlášení vázaná na konkrétní řádek a pravidlo.
Výstup \texttt{cppcheck} ve tvaru \texttt{[file.cpp:47]: (error) Memory leak: p} je precizní, avšak pro studenta, který nemusí znát základy správy paměti, je pouze obtížně pochopitelný bod bez vysvětlení.
LLM je schopen generovat vysvětlení přirozeným jazykem, které studentovi popíše, v čem spočívá problém, proč je nebezpečný a jak ho opravit.
Souhrnné porovnání obou přístupů je uvedeno v \tabref[tabulce]{tab:llm-vs-static-summary}.

\begin{table}[H]
    \centering
    \resizebox{\textwidth}{!}{%
        \begin{tabular}{lll}
            \toprule
            \textbf{Vlastnost}                   & \textbf{Statický analyzátor} & \textbf{Velký jazykový model} \\
            \midrule
            Deterministický výstupu              & Ano                          & Ne (pravděpodobnostní)        \\
            Formální verifikace pravidel         & Ano                          & Ne                            \\
            Analýza řízení toku (CFG)            & Ano                          & Omezeně / heuristicky         \\
            Detekce syntaktických chyb           & Ano (spolehlivě)             & Omezeně                       \\
            Práce s kontextem zadání úlohy       & Ne                           & Ano                           \\
            Hodnocení vhodnosti algoritmu        & Ne                           & Ano (při znalosti zadání)     \\
            Hodnocení čitelnosti a stylu         & Omezeně (dle pravidel)       & Ano                           \\
            Vysvětlení přirozeným jazykem        & Ne (technická hlášení)       & Ano                           \\
            Škálovatelnost na libovolný kód      & Ano                          & Limitováno kontextem          \\
            Detekce bezpečnostních zranitelností & Ano (definované vzory)       & Omezeně                       \\
            \bottomrule
        \end{tabular}
    }
    \caption{Porovnání statických analyzátorů a velkých jazykových modelů při analýze zdrojového kódu.}
    \label{tab:llm-vs-static-summary}
\end{table}

\endinput